% !TeX root = ../papers/nc-p3.tex
\section{中心の周期四の類と固定Brauer類における窓の一意性}
\label{sec:thirty-node-brauer-windows}
\begin{theorem}\label{thm:fixed-brauer-center-windows}
前節の一般中心 $X$、その滑らかな開集合 $U$、窓から来るBrauer類 $\alpha$ に対し、
$\alpha$ の周期と指数はともに4である。
中心を $X$ と同定し同じ類 $\alpha$ を持つ、この三十節点の開集合上の窓は、
添字移動を許せば一つである。中心の同定を保つ添字保持同型では、ちょうど四つの連続移動になる。
\end{theorem}
\begin{proof}[因子と局所次数]
$V$ を $\mathbb P^4$ の三十点での吹上げ、$Y$ を $X$ の狭義変換とする。
$H$ を超平面、$E$ を例外因子の和とし、$Y\sim4H-2E$、$K_V=-5H+3E$ である。
$Y$ は滑らかで、$\operatorname{Pic}(U)\otimes\mathbb Q$ は $H$ で張られることを示す。
標数零の有限型データの基礎体を複素数に移してコホモロジーを計算してよい。
$\mathcal O_V(-Y)=\mathcal O_V(-4H+2E)$ のコホモロジーは全て零である。
実際、引き戻した $\mathcal O(-4)$ から例外因子を一枚ずつ足すとき、商は
$\mathbb P^3$ 上の $\mathcal O(-1),\mathcal O(-2)$ である。
特に $H^1(Y,\mathcal O_Y)=H^2(Y,\mathcal O_Y)=0$ となる。

$H^2(Y,\mathcal O_Y(-Y))$ の双対は $H^1(Y,\mathcal O_Y(3H-E))$ である。
$V$ 上では $3H-E$ の第一コホモロジーは三十点への三次式の全射により零で、
制限列の核 $-H+E$ も非コホモロジー的なので、この双対も零である。
また $H^2(V,\Omega_V^1(-Y))=0$ である。吹上げの微分の余核は
各例外 $\mathbb P^3$ の $\Omega^1$ の直和である（チャート $x_1=t,x_j=tu_j$ で直ちに分かる）。
この商を $-Y$ 捩ると $\Omega^1_{\mathbb P^3}(-2)$ で、Euler列から非コホモロジー的である。
引き戻しの部分は $\Omega^1_{\mathbb P^4}(-4)$ の引き戻しに $2E$ を足したもので、
先と同じ理由でコホモロジーは元の束と同じ、特に第二は零である。
よって余法線列と制限列から $H^1(V,\Omega^1_V)\to H^1(Y,\Omega^1_Y)$ は全射となる。
左辺は $H$ と各例外因子の類で張られる（同じ微分列、または滑らかな点の吹上げの公式）。
Hodge分解と第一Chern類から $\operatorname{Pic}(Y)\otimes\mathbb Q$ はこれらで張られる。
滑らかな開集合の因子を閉包に取れるので $\operatorname{Pic}(Y)\to\operatorname{Pic}(U)$ は全射であり、
主張を得る。従って各完備節点の穴あきスペクトル $U_p$ への
$\operatorname{Pic}(U)\to\operatorname{Pic}(U_p)$ は零である。
実際右辺は $\mathbb Z$ で、二次曲面の二罫線の差による次数で測られる。
この局所計算は、節点の吹上げ上で除数を延長し、
例外 $Q=\mathbb P^1\times\mathbb P^1$ の $\mathcal O(-1,-1)$ を割っても得られる。
その吹上げのPicard群は $\operatorname{Pic}(Q)$ と一致する。
二罫線の束は行列の像・核の二射影から延び、厚み上での一意性は
$H^1(Q,\mathcal O(n,n))=0$、完備化との比較は固有形式存在定理による。

$U$ 上で分裂のgerbeを取り、窓の四冪等元をその上の重み1の階数4束 $E_i$ とする。
この規約で共通のねじれ類を $\alpha$ と書く。
$4\alpha=0$ は階数4の行列代数から従う。
もし $2\alpha=0$ なら、このgerbeに重み2の直線束 $P$ がある
（分裂の遷移の二乗の障害が $2\alpha$ である）。
$\det E_i\otimes P^{-2}$ は $U$ 上の通常の直線束である。
既知の二十節点のどれかで局所gerbeを自明化する。
原始冪等元の階数と各単純の頂点次元1,3により、そこでの $E_i$ は共通の直線捩りを除き
$\mathcal O\oplus I^{\oplus3}$ または $\mathcal O^{\oplus3}\oplus I$ である。
これは $R^4\oplus I^{12}$ 上の $i$ 番目の冪等元の像を取ったものである。
従ってその行列式は局所類群で奇数、$P^{-2}$ の寄与は偶数で、先の制限零に反する。
よって周期は4である。指数4は系\ref{cor:quartic-center-unramified-brauer}による。
\end{proof}

\begin{proof}[一般超平面の定義体での復元]
中心とBrauer類を同定した二窓の四束列を同じgerbe上で $E_i,E'_i$ とする。
$E_{i+4}=E_i(1)$ として全整数に延長する。最初の指数は、四連続の大切断の
順方向射が $\mathcal Hom(E_i,E_j)$ の切断になる規約で選ぶ。
$K$ を超平面パラメータの関数体、$S$ をその上の一般切断とする。
これは $U$ に含まれる滑らかな四次K3である。普遍超平面は $U$ 上では
射影三次元空間束なので、$K(S)$ は $k(X)$ の純超越拡大である。
従って $K(S)$ でも $\alpha$ の指数は4である。特に $K$ 上の各 $E_i|_S,E'_i|_S$ は
より小さい正階数のねじれ部分層を持たず、傾き安定である。以下のHomと双対性は
全て $K$ 上で計算し、幾何学的一般超平面でのBrauer群の単射を必要としない。

同階数の二束の行列式の比は通常の直線束であり、
$U$ から来るので $S$ 上の数値類は整数倍の $h=H|_S$ である。
実際有理倍であることは前段による。幾何的K3で $h^2=4$ かつ格子が偶であるため
$h$ は原始で、Picard群に捩れもない。
$d=j-i$ とし、$E_i,E_j$ のHom束の $mh$ 捩りのEuler多項式は
\[
 {4m+d+3\choose3}-{4m+d-1\choose3}
   =32m^2+16dm+2d^2+2 .
\]
中心超平面で割る切断列と元の標準Hilbert関数を使った。
従って $c_1(\det E_j/\det E_i)h=4d$、
その類は $dh$ である。$E'_i$ にも同じ式が成り立つ。
よって一つの整数 $s$ を取れば全ての $i$ で
$\det E_{i+s}|_S$ と $\det E'_i|_S$ は同じ数値類を持つ。

この二束を $L,M$ と略記する。上式の $d=0,m=0$ により
$\chi(\mathcal End L)=\chi(\mathcal End M)=2$ である。
通常の十六次元束 $D=\mathcal Hom(L,M)$ は数値的第一Chern類が零で、
$\mathcal End D=\mathcal End L\otimes\mathcal End M$ である。
第二Chern指標で比較すると $2\operatorname{ch}_2D=
\operatorname{ch}_2(\mathcal End L)+\operatorname{ch}_2(\mathcal End M)$。
K3上のRiemann--Rochで $\chi(D)=2$ となり、Serre双対から一方向に非零射がある。
一般点の階数は4の倍数なのでこの射は単射で、同じ行列式次数から余核に因子台はない。
行列式は零因子を持たない束の射となるので全点で同型である。従って $E_{i+s}|_S\simeq E'_i|_S$。

四束の同型を周期的に捩って全位置で取る。各隅の切断列で中間コホモロジーが消えるので
経路長が1と2の切断とそれらの積は制限後も全く同じである（核は経路長から4を引いた空間）。
よって四文字と全関係の同型を $K$ 上で得る。これは標準的な周期を備えた四窓に使える。
実際四次socle直線から得る周期表示では三文字と二関係の同型の第四文字への延長は
スカラーを除き一意であり、上の周期的な束同型は全四関係を保っている。
その中心の射影自己同型は定義体 $k$ 上の有限群の元である。
四次中心の制限では超平面方程式だけが零となる。束の同型は各長さ4の積も制限上で運ぶので、
この自己同型は一般超平面を保つ。$k$ 上の射影変換がその一般超平面を固定すれば恒等である。
中心の同定を保つ窓同型の方程式は有限型であるため $K$ 上の解から代数閉体 $k$ 上の解が従う。

逆に中心を固定した同型は、$U$ 上の分裂束の列を共通の通常の直線捩りまで運ぶ。
その行列式の差は $S$ 上で $4h$ の整数倍である。従って移動数は法4で一意に決まり、
四つの移動は互いに異なる。各移動は元の周期的な窓の開始頂点を変えたもので、
同じ中心と同じBrauer類を持つ。
\end{proof}
